% Modernised reading — fonds Alexandre Grothendieck, shelfmark 89 (pages 1-20).
%
% This is an INTERPRETATION, not a transcription. It re-reads the pages in
% today's notation and vocabulary, states the hypotheses the manuscript leaves
% implicit, and drops the critical apparatus so that the mathematics reads
% straight through. Everything the brackets and underlines carried moves into
% a footnote. It is held to being mathematically correct as it stands; where
% the manuscript is loose or mistaken the true statement appears and the
% footnote says what the page has. In any dispute of substance the
% transcription governs: whoever cites Grothendieck cites batch-01.fr.tex.
%
% Scope: a single document for the whole folder. The shelfmark holds one
% batch (pages 1-20), read whole before any of this was written. Page 1 is a
% cover carrying « les ensembles 3-3 ».
%
% Spine. Four runs: I (2-8) hexagons and the automorphisms of a 3-3 set with
% an identification of orientations; II (9-14) planes over F_3, their pairs
% of lines, combinatorial squares; pages 15-18 the same reconstruction
% question for a free module of rank 2 over a ring; pages 19-20 combinatorial
% polygons and orientations. The connection that makes I a special case of II
% is THIS READING'S: a 3-3 set S = S_1 u S_2 canonically gives the plane
% V_S = D_1 (+) D_2 with D_i = omega_{S_i} u {0}, and the torsor P_S = S_1 x
% S_2; the six bijections S_1 -> S_2 are the affine lines of P_S in the two
% remaining directions, so the S' of pages 7-8 is the underlying 3-3 set of
% the equivalence (8) of page 12. The page never says this.
%
% Notation fixed for the whole folder. Ens_n is the groupoid of sets with n
% elements. For a three-element set E, omega_E is the two-element set of its
% cyclic orders, a {+-1}-torsor. A (3,3)-set (the pages' « ensemble 3-3 ») is a
% six-element set S with a partition into two three-element classes S_i,
% indexed by the two-element set T. The contracted product of torsors is
% written ∧, as on the pages; Gamma_omega = Z/3 ∧_{+-1} omega. For a line D
% over F_3 (or a rank-one free module), D^* is the set of its generators.
% Page 3's G is kept; the page's G' names two different groups, set apart
% here as K (the kernel of G -> S_T, order 18, page 3) and N (the subgroup
% generated by the antipodisms, order 6, pages 4-6). The sans-serif 𝖳 of
% page 6 (a Gamma_omega-torsor) is kept apart from the index set T.
%
% Substantive departures from the page, with pages.
%  - Pages 4-5: the subgroup generated by the three antipodisms is written G'
%    (and first V, « 2-groupe (invariant) »); it has order 6, is isomorphic to
%    S_3, and is not the G' of page 3. Renamed N.
%  - Pages 5-6: the page asks « Quel est le groupe quotient G/G' ? » and
%    breaks off before answering; this reading answers G = N x M ≅ S_I x S_𝖳,
%    with M generated by the three orientation-reversing involutions of page
%    6, and says the answer is its own.
%  - Page 10: « partitions de type (3,3) » for the second pair is (3,3,3).
%  - Page 14: the numeral (12) is written as (11).
%  - Pages 15-18: the equivalence u ∧ v = v ∧ w = w ∧ u <=> u + v + w = 0
%    holds over any ring, not only a field; the reconstruction is stated, as
%    on the page, for a field.
%  - Page 20: the theorem needs Pi connected of order >= 3; c) is false for
%    the square with the condition « f(s) adjacent to s » alone (two
%    reflections also satisfy it) and needs the margin's f(a) ≠ a; for the
%    digon the map omega |-> f_omega is not injective.
%
% Announced and never established in the folder: the sentence of page 8 on a
% hexagon « and a vertex on it » (breaks off); the answer to G/G' (page 5);
% the list of properties needed on pages 18 (only a) is written); the
% reconstruction of V over a ring rather than a field (page 17); the
% « Foncteur réalisation topologique » and the relation with topological
% orientations (pages 19-20), named only.
%
% Pass: Opus 5 (claude-opus-5), 2026-09-13 — first pass, unchecked against the pages by a human.
\documentclass[11pt,a4paper]{article}
% Le préambule commun à toutes les transcriptions du fonds Grothendieck.
%
% Il tient en une idée : **l'incertitude se compose, elle ne se résout pas.**
% Une transcription qui lisse un mot douteux a détruit la seule chose pour
% laquelle on la faisait — un lecteur doit pouvoir savoir, à chaque endroit,
% ce qui était sur le feuillet et ce qui a été ajouté. Les six macros qui
% suivent sont donc l'appareil critique, et non de la décoration.
%
% Les mêmes commandes sont reconnues par `scripts/render.mjs`, qui produit la
% vue de lecture du site. Une macro ajoutée ici doit l'être aussi là-bas,
% faute de quoi le rendu échouera bruyamment — ce qui est voulu : un rendu
% silencieusement faux est pire qu'un rendu absent.

\ProvidesPackage{grothendieck}[2026/08/08 Transcriptions du fonds Grothendieck]

\usepackage{amsmath,amssymb,amsthm}
\usepackage{mathrsfs}
\usepackage{stmaryrd}
% Les diagrammes commutatifs sont partout dans ce fonds : c'est la langue
% même de la géométrie algébrique de Grothendieck, et une transcription qui
% les rendrait en prose ne transcrirait rien.
\usepackage{tikz-cd}
% Français par défaut — c'est la langue des feuillets — mais l'anglais est
% chargé aussi : la traduction et le résumé sont des documents anglais, et la
% typographie française (espaces avant les deux-points) y serait une faute.
% Ces documents basculent par \selectlanguage{english} en tête de corps.
\usepackage[english,french]{babel}
\usepackage{fontspec}
\usepackage{geometry}
\geometry{a4paper,margin=25mm,marginparwidth=28mm}
\usepackage{marginnote}
\usepackage{xcolor}
% ulem, not soul. `soul`'s \st and \ul are built on a character-by-character
% scan that XeLaTeX's font machinery breaks: the document fails with an
% undefined control sequence rather than degrading. ulem does the same two jobs
% and is Unicode-engine safe. `normalem` stops it hijacking \emph, which the
% transcriptions use for Grothendieck's own emphasis.
\usepackage[normalem]{ulem}
\usepackage{hyperref}
\hypersetup{colorlinks=true,linkcolor=black,urlcolor=black,pdfborder={0 0 0}}

% --- Métadonnées du lot -----------------------------------------------------
% Reprises telles quelles par le script de rendu, qui les affiche en tête de
% la vue de lecture. Elles fixent ce que le fichier prétend couvrir : sans
% elles, une transcription flotte sans point d'attache dans un fonds de seize
% mille feuillets.

\newcommand{\folder}[1]{\gdef\grothFolder{#1}}
\newcommand{\batch}[1]{\gdef\grothBatch{#1}}
\newcommand{\leaves}[2]{\gdef\grothFirst{#1}\gdef\grothLast{#2}}
\newcommand{\foldertitle}[1]{\gdef\grothFoldertitle{#1}}
\gdef\grothFolder{?}\gdef\grothBatch{?}\gdef\grothFirst{?}\gdef\grothLast{?}\gdef\grothFoldertitle{}

% --- Le résumé --------------------------------------------------------------
%
% Il ouvre la lecture modernisée, sous le titre « Résumé », et il est composé
% en petit. Ce n'est pas un ornement : c'est le seul passage du document qui
% ne soit pas des mathématiques, et un lecteur doit pouvoir le reconnaître
% avant de l'avoir lu — puis le sauter s'il connaît déjà le sujet, ou s'y
% arrêter s'il ne le connaît pas. Le filet qui le suit marque la reprise du
% corps.

\newenvironment{resume}
  {\small\setlength{\parskip}{0.45em}}
  {\par\medskip\hrule\medskip}

% --- Mots-clés ---------------------------------------------------------------
%
% En anglais, en clôture du résumé de la lecture modernisée : le vocabulaire
% moderne sous lequel chercher ce que les feuillets construisent. Le manifeste
% du site (`npm run manifest`) les extrait pour étiqueter la cote entière —
% cette ligne est la source unique des tags, il n'y a pas de fichier à côté.
\newcommand{\keywords}[1]{\par\smallskip\noindent
  {\footnotesize\textsc{Keywords} --- \textit{#1}}\par}

% --- L'appareil critique ----------------------------------------------------

% Le numéro de feuillet, en marge. C'est l'ancre qui permet de régler un
% désaccord sur une formule en regardant *un* feuillet plutôt qu'une chemise
% de six cents. Le site s'en sert aussi pour tourner les pages du fac-similé
% quand on fait défiler la transcription.
\newcommand{\leaf}[1]{%
  \marginnote{\footnotesize\color{gray!70}\textsf{#1}}%
  \label{leaf:#1}\ignorespaces}

% Illisible : on ne devine pas. Un mot inventé qui se lit comme les autres est
% le pire résultat possible.
\newcommand{\ill}{\textcolor{red!60!black}{[\,\ldots\,]}}

% Lecture douteuse : le mot est donné, et signalé comme incertain.
\newcommand{\uncertain}[1]{\uline{#1}}

% Ajout de l'éditeur — une lettre manquante, un mot sous-entendu.
\newcommand{\add}[1]{\textcolor{blue!50!black}{[#1]}}

% Biffé par Grothendieck lui-même. Ce qu'il a rayé fait partie du document :
% une rature dit souvent où la pensée a changé de direction.
\newcommand{\struck}[1]{\sout{#1}}

% Note du transcripteur, sur l'état matériel du feuillet.
\newcommand{\note}[1]{\par\smallskip{\small\color{gray!60!black}\textit{[#1]}}\par\smallskip}

% Note marginale de Grothendieck, distincte du corps du texte.
\newcommand{\marginal}[1]{\par\smallskip\noindent\hspace{1em}%
  {\small\color{gray!60!black}#1}\par\smallskip}

% --- Titre ------------------------------------------------------------------

\AtBeginDocument{%
  \begin{center}
    {\large\bfseries Fonds Alexandre Grothendieck --- cote n\textsuperscript{o}~\grothFolder}\\[2pt]
    {\itshape\grothFoldertitle}\\[4pt]
    {\small feuillets \grothFirst--\grothLast\ (lot \grothBatch)}\\[2pt]
    {\footnotesize\color{gray!60!black}Transcription établie d'après le fac-similé numérisé
     par l'Université de Montpellier.\\
     Les lectures incertaines sont soulignées, les illisibles notées [\,\ldots\,],
     les ajouts entre crochets.}
  \end{center}
  \vspace{1em}}

\endinput


\folder{89}
\pages{1}{20}
\dating{s.d. — le groupe « Géométrie et topologie combinatoire » (69 à 102) est daté 1976-[vers 1986]}
\watermark{Édition de démonstration\\interprétation personnelle de l'œuvre}
\foldertitle{[Hexagone] combinatoire. [Ensembles 3-3] : notes manuscrites (s.d.) — lecture modernisée du dossier entier}

\begin{document}

\section*{Résumé}

\begin{resume}
Prenez six objets et rangez-les en deux groupes de trois. C'est la structure
la plus pauvre qui soit, et ces vingt pages montrent qu'elle en cache
plusieurs autres, plus riches, qu'on peut y lire sans rien ajouter — ou
presque rien.

Première lecture : un hexagone. Les six sommets d'un hexagone régulier se
répartissent en deux triangles inscrits, et chaque sommet est déterminé par
son triangle et par la grande diagonale qui le traverse. Grothendieck
remarque que se donner un hexagone, à la façon d'un combinatoricien — six
points, six arêtes, sans règle ni compas —, revient exactement à se donner un
ensemble à trois éléments (les diagonales) et un ensemble à deux éléments (les
triangles). Il étudie ensuite les symétries d'un « ensemble 3-3 » dont les deux
triples ont été orientés de façon compatible : ce groupe a trente-six
éléments, il contient trois « antipodies » qui engendrent un groupe de
permutations de trois objets, et la question de ce qui reste quand on les
quotiente est posée et laissée ouverte sur la page.

Deuxième lecture : une géométrie. Sur le corps à trois éléments, un plan a
neuf points et quatre directions de droites. Choisir deux de ces directions,
c'est découper les neuf points en un quadrillage $3 \times 3$ ; les deux autres
directions donnent un second quadrillage. Grothendieck en tire un va-et-vient
entre deux ensembles 3-3, et une correspondance exacte entre les plans sur
$\mathbf{F}_3$ munis de deux directions et les carrés combinatoires, où le carré
dual — dont les sommets sont les arêtes du premier — correspond aux deux
autres directions. La page 14 aboutit à une relation d'une étrange précision :
passer du carré à son dual puis revenir donne l'antipodie, $-1$, et cela tient
à ce que $1 + 1 = -1$ dans ce corps.

Les deux lectures sont en fait la même : les six bijections entre les deux
triples d'un ensemble 3-3 sont exactement les droites du plan à neuf points
dans les deux directions restantes. Les pages ne le disent pas, mais c'est ce
qui fait de ce dossier un tout. Il se prolonge dans deux directions. Les pages
15 à 18 demandent comment reconstituer un plan vectoriel sur un anneau
quelconque à partir de ses seules droites et de la façon dont deux droites se
complètent, et découvrent qu'une somme $u + v + w = 0$ se lit sur le
déterminant. Les pages 19 et 20, d'une autre main plus rapide, posent les
définitions des polygones combinatoires et énoncent qu'un polygone connexe a
exactement deux orientations.

Ce qu'on y voit de sa manière : il ne fixe jamais un choix quand il peut garder
tous les choix ensemble. Là où l'on dirait « orientons l'hexagone », il écrit
le groupe des rotations « tordu » par l'ensemble des deux orientations, et
toute la suite consiste à faire apparaître ce qui ne dépend d'aucun choix.

\keywords{combinatorial hexagon, torsor, contracted product, wreath product, Coxeter presentation, affine plane over F3, dual square, projective line over a ring, unimodular vector, determinant, combinatorial polygon, orientation of a cycle graph}
\end{resume}

\pagerange{2}{20}
\section*{Le fil du dossier, et les conventions}

Le feuillet de couverture (page 1) porte seulement « les ensembles 3-3 ». Il n'y
a pas de pagination de sa main ; deux chiffres romains, I (page 2) et II
(page 9), ouvrent les deux premiers ensembles.

\begin{itemize}
  \item \textbf{Pages 2 à 8 (I) : hexagones et ensembles 3-3.} Un hexagone
    combinatoire équivaut à un couple (diagonales, triangles) ; le groupe
    d'un ensemble 3-3 muni d'une identification de ses orientations, ses
    antipodies, le torseur $\mathsf{T}$ ; puis, pages 7 et 8, l'ensemble $S'$ des
    six scindages, lui-même un ensemble 3-3. La série s'interrompt au milieu
    d'une phrase.
  \item \textbf{Pages 9 à 14 (II) : plans sur $\mathbf{F}_3$ et carrés.} Paires de
    droites, décompositions d'un plan affine à neuf points, l'équivalence (8)
    entre ensembles 3-3 indexés par deux paires complémentaires, les carrés
    combinatoires et leur dual, la relation (13).
  \item \textbf{Pages 15 à 18 : modules libres de rang $2$.} Reconstituer $V$ à
    partir de ses droites, des torseurs de générateurs et du déterminant. La
    liste des propriétés nécessaires s'arrête après a).
  \item \textbf{Pages 19 et 20 : polygones combinatoires.} Trois points de vue,
    orientations, un théorème en trois parties. D'une autre main, plus rapide.
\end{itemize}

\textbf{Le lien entre I et II}, qui est de cette lecture. Un ensemble 3-3
$S = S_1 \sqcup S_2$ définit canoniquement deux droites sur $\mathbf{F}_3$,
$D_i = \omega_{S_i} \sqcup \{0\}$ (page 13 : une droite sur $\mathbf{F}_3$ est
reconstituée par ses deux générateurs), un plan $V_S = D_1 \oplus D_2$ et un
$V_S$-torseur $P_S = S_1 \times S_2$ à neuf points. Les deux autres droites de
$V_S$ sont les graphes des deux isomorphismes $D_1 \to D_2$, et comme toute
bijection entre deux ensembles à trois éléments est affine, les six bijections
$S_1 \to S_2$ sont exactement les six droites affines de $P_S$ dans ces deux
directions. L'ensemble $S'$ des pages 7 et 8 est donc l'ensemble 3-3 que
l'équivalence (8) de la page 12 associe à $S$, et l'identification (A) des
orientations de la page 3 revient au choix de l'une de ces deux directions.

\textbf{Conventions, valables pour tout le dossier.} $\mathrm{Ens}_n$ est le
groupoïde des ensembles à $n$ éléments. Pour un ensemble $E$ à trois éléments,
$\omega_E$ est l'ensemble de ses deux ordres circulaires, torseur sous
$\{\pm 1\}$ par inversion. Un \emph{ensemble 3-3} est un ensemble $S$ à six
éléments muni d'une partition en deux classes $S_i$ à trois éléments, indexées
par un ensemble $T$ à deux éléments. Le produit contracté de torseurs est noté
$\wedge$, comme sur les pages. Pour un groupe abélien $A$ sur lequel $\{\pm 1\}$
agit par $x \mapsto -x$, et un $\{\pm 1\}$-torseur $\omega$, on note
$A \wedge \omega = (A \times \omega)/\bigl((x, o) \sim (-x, -o)\bigr)$ la forme
tordue de $A$ par $\omega$ ; ainsi $\Gamma_\omega = \mathbf{Z}/3\mathbf{Z} \wedge \omega$.
Pour une droite $D$ sur $\mathbf{F}_3$, ou un module libre de rang $1$, $D^{*}$
désigne l'ensemble de ses générateurs. Le groupe noté $G'$ par les pages
désigne deux groupes différents ; on les nomme ici $K$ (page 3) et $N$ (pages 4
à 6).

\pagerange{2}{8}
\section*{Hexagones et ensembles 3-3 (pages 2 à 8)}

\pagerange{2}{3}
\subsection*{Un hexagone équivaut à trois diagonales et deux triangles}

Un \emph{hexagone combinatoire} $H$ est un polygone combinatoire à six sommets,
au sens des pages 19 et 20 : un ensemble $S$ de six sommets et un ensemble $A$
d'arêtes, chaque sommet sur exactement deux arêtes, formant un seul cycle. Ses
rotations forment un groupe cyclique d'ordre $6$, qui agit simplement
transitivement sur $S$ ; mais l'identifier à $\mathbf{Z}/6\mathbf{Z}$ demande de
choisir un sens. Grothendieck évite le choix : si $\omega_H$ est l'ensemble des
deux orientations de $H$, le groupe des rotations est canoniquement
$\mathbf{Z}/6\mathbf{Z} \wedge \omega_H$, et $S$ en est un torseur. Comme
\[
\mathbf{Z}/6\mathbf{Z} \wedge \omega_H \simeq (\mathbf{Z}/3\mathbf{Z} \wedge \omega_H) \times
(\mathbf{Z}/2\mathbf{Z} \wedge \omega_H), \qquad \mathbf{Z}/2\mathbf{Z} \wedge \omega_H = \mathbf{Z}/2\mathbf{Z},
\]
un torseur sous ce groupe est un couple de torseurs :
\begin{itemize}
  \item le quotient $\Delta = S/(\mathbf{Z}/2\mathbf{Z})$, ensemble des trois paires de
    sommets opposés, c'est-à-dire des trois grandes diagonales, torseur sous
    $\mathbf{Z}/3\mathbf{Z} \wedge \omega_H$ — ce qui revient à la donnée d'une
    bijection $\omega_\Delta \simeq \omega_H$ ;
  \item le quotient $T = S/(\mathbf{Z}/3\mathbf{Z} \wedge \omega_H)$, ensemble à deux
    éléments des triangles inscrits.
\end{itemize}
Comme $\omega_H$ se retrouve à partir de $\Delta$, la donnée de $H$ équivaut à
celle du couple $(\Delta, T)$, et l'on a une bijection canonique
\[
S \simeq \Delta \times T :
\]
un sommet est connu dès qu'on connaît le triangle et la diagonale qui le
contiennent, et ces deux données sont arbitraires\footnote{Le mot qui désigne
le sommet est illisible sur la page 2 ; « sommet » est ce que l'argument
demande. Le « au même » qui précède la paire de torseurs est une lecture
incertaine.}. Réciproquement, sur $\Delta \times T$, on reconstitue l'hexagone en
joignant $(d, i)$ à $(d', i')$ si et seulement si $d \neq d'$ et $i \neq i'$ :
c'est le graphe biparti complet $K_{3,3}$ privé des trois arêtes $\{(d, i), (d,
i')\}$, qui est un cycle de longueur $6$\footnote{Cette reconstruction explicite
est de l'édition ; la page dit seulement « on récupère $H$ par un iso.
can. ».}. D'où l'énoncé que la page porte en marge :
\[
\{\text{hexagones combinatoires}\} \simeq \mathrm{Ens}_3 \times \mathrm{Ens}_2 ,
\]
équivalence de groupoïdes ; on vérifie les groupes d'automorphismes, le groupe
diédral d'ordre $12$ étant isomorphe à $\mathfrak{S}_3 \times \mathbf{Z}/2\mathbf{Z}$.

Cette décomposition fait de $S$ un ensemble 3-3, de classes les deux triangles
$S_i = \Delta \times \{i\}$, $i \in T$. Chaque $S_i$ est en bijection avec $\Delta$,
d'où $\omega_{S_i} \simeq \omega_\Delta$ pour tout $i$, et donc un isomorphisme
canonique
\[
\bigwedge_{i \in T} \omega_{S_i} \simeq \mathbf{1},
\]
le $\{\pm 1\}$-torseur trivial : le produit contracté de deux copies d'un même
torseur sous $\{\pm 1\}$ est canoniquement trivial.

\pagerange{3}{6}
\subsection*{Le groupe d'un ensemble 3-3 muni d'une identification des orientations}

Considérons maintenant, sur un ensemble $S$ à six éléments, la structure formée
d'une partition de type $(3,3)$ indexée par $T$ et d'un isomorphisme
\[
(\mathrm{A}) \qquad \bigwedge_{i \in T} \omega_{S_i} \simeq \mathbf{1},
\]
c'est-à-dire d'une identification des deux ensembles d'orientations
$\omega_{S_1} \simeq \omega_{S_2}$\footnote{La page parle d'un « système
transitif d'isomorphismes » entre les $\omega_{S_i}$ ; pour deux ensembles, c'est
un isomorphisme. Elle appelle ensuite cette structure, semble-t-il,
« unisignée », lecture incertaine.}. Son groupe d'automorphismes $G$ est
d'ordre $36$. En effet, $G$ se surjecte sur $\mathfrak{S}_T \simeq \mathbf{Z}/2\mathbf{Z}$, et son
noyau
\[
K = \ker\Bigl(\prod_{i \in T} \mathfrak{S}_{S_i} \xrightarrow{\ \prod_i \mathrm{sg}_i\ } \mathbf{Z}/2\mathbf{Z}\Bigr)
\]
est formé des couples de permutations de même signature, d'ordre $18$\footnote{La
page note ce noyau $G'$, dessine la suite exacte $1 \to G' \to G \to \mathfrak{S}_T \to 1$
et, verticalement, $G' \to \prod_i \mathfrak{S}_{S_i} \to \mathbf{Z}/2\mathbf{Z} \to 1$.}.
Un élément qui échange les deux classes par $a : S_1 \to S_2$ et
$b : S_2 \to S_1$ préserve (A) si et seulement si $b\,a$ est paire.

\textbf{Hexagones subordonnés et antipodies.} Un hexagone sur $S$ dont les
triangles sont les $S_i$ est déterminé par l'appariement de ses sommets opposés,
c'est-à-dire par une bijection $u : S_1 \to S_2$, et l'identification induite des
orientations est celle que transporte $u$. L'hexagone est donc \emph{subordonné}
à la structure (A) exactement quand $u$ transporte l'orientation selon (A) :
parmi les six bijections $S_1 \to S_2$, trois conviennent. À chacune correspond
l'\emph{antipodie} $\sigma_H$, involution de $S$ égale à $u$ sur $S_1$ et à
$u^{-1}$ sur $S_2$ : c'est un élément de $G$ au-dessus de la transposition de $T$,
et il détermine $H$\footnote{La page dit « inverseurs d'orientations », lecture
incertaine, là où l'argument demande que $u$ \emph{respecte} l'identification
(A). Une note marginale en éventail sur la même question ne se lit que par
fragments.}. Le groupe des automorphismes de $H$ est le centralisateur
$Z(\sigma_H)$, d'ordre $12$ et d'indice $3$ dans $G$ ; il n'est pas distingué, il
est son propre normalisateur, et ses trois conjugués sont les centralisateurs des
trois antipodies, qui forment une classe de conjugaison de $G$ : les trois
hexagones subordonnés sont échangés par $G$\footnote{« son propre normalisateur »
est notre lecture d'une insertion incertaine, « et en particulier le groupe
normalisateur », rattachée au mot « distingué ».}.

\textbf{Le groupe engendré par les antipodies.} Soit $I = \{i, j, k\}$ l'ensemble
des trois antipodies $\sigma_i$ et $N$ le sous-groupe de $G$ qu'elles engendrent,
distingué puisqu'engendré par une classe de conjugaison\footnote{La page le
nomme d'abord $V$, « $2$-groupe (invariant) », puis $G'$. Ce n'est ni un
$2$-groupe — il est d'ordre $6$ — ni le $G'$ de la page 3, d'où le nom $N$.}. On
vérifie que $\sigma_k = \sigma_i \sigma_j \sigma_i$, et par permutation circulaire
$\sigma_i = \sigma_j\sigma_k\sigma_j$, d'où $(\sigma_i\sigma_j)^3 = 1$. L'action de $N$ par
conjugaison sur $I$ donne un isomorphisme
\[
N \xrightarrow{\ \sim\ } \mathfrak{S}_I ,
\]
qui envoie $\sigma_i$ sur la transposition de $I$ fixant $i$. Ainsi $N$ est
présenté par les générateurs $\sigma_i$ ($i \in I$) et les $3 + 6 = 9$ relations
$\sigma_i^2 = 1$, $\sigma_k = \sigma_i\sigma_j\sigma_i$ ; ou encore, pour $i \neq j$ fixés,
$\sigma = \sigma_i$, $\sigma' = \sigma_j$, par
\[
\sigma^2 = \sigma'^2 = (\sigma\sigma')^3 = 1 ,
\]
la présentation de Coxeter de $\mathfrak{S}_3$.

\textbf{Le quotient $G/N$.} La page demande ce qu'est ce quotient, d'ordre $6$.
La structure (A) équivaut à la donnée d'un ensemble $\omega$ à deux éléments —
l'orientation commune — et de deux torseurs $S_i$ sous $\Gamma_\omega =
\mathbf{Z}/3\mathbf{Z} \wedge \omega$. On forme le $\Gamma_\omega$-torseur
\[
\mathsf{T} = \bigwedge_{i \in T} S_i ,
\]
et $G$ agit sur le couple $(\Gamma_\omega, \mathsf{T})$. \emph{Le sous-groupe $N$ y agit
trivialement.} En effet, une antipodie $\sigma$ donnée par $u : S_i \to S_j$ induit
$x \wedge y \mapsto u^{-1}y \wedge ux$ ; comme $u$ respecte l'orientation, il commute
à $\Gamma_\omega$, et, pour $a \in S_i$, $b = ua$, $x = \gamma a$, $y = \gamma' b$ :
\[
\sigma(\gamma a \wedge \gamma' b) = \gamma' a \wedge \gamma b = \gamma'\gamma\,(a \wedge b)
= \gamma\gamma'\,(a \wedge b) = \gamma a \wedge \gamma' b ,
\]
$\Gamma_\omega$ étant commutatif. Donc $G/N$ agit sur $(\Gamma_\omega, \mathsf{T})$.

La page indique ensuite que $\mathsf{T}$ se réalise canoniquement comme l'ensemble des
trois involutions de $S$ qui échangent les deux classes \emph{et} renversent
l'orientation commune, et que ces involutions se conjuguent entre elles. On peut
achever\footnote{La réponse qui suit n'est pas sur la page, qui s'arrête sur
« (en fait … l'a déjà fait …) ». Elle se vérifie sur le modèle
$S = \mathbf{Z}/3\mathbf{Z} \times \{1, 2\}$, où les antipodies sont
$(d, 1) \mapsto (d + c, 2)$ et les involutions renversantes $(d, 1) \mapsto (s - d, 2)$.} :
ces trois involutions engendrent un sous-groupe $M \simeq \mathfrak{S}_3$, qui commute à
$N$ et le rencontre trivialement, de sorte que
\[
G = N \times M \simeq \mathfrak{S}_I \times \mathfrak{S}_{\mathsf{T}}, \qquad G/N \simeq M \simeq \mathfrak{S}_{\mathsf{T}} ,
\]
$M$ agissant fidèlement sur $\mathsf{T}$. Le groupe $G$ est donc le produit de deux
groupes symétriques, chacun engendré par trois involutions qui échangent les
triangles — celles qui respectent l'orientation et celles qui la renversent.

\pagerange{7}{8}
\subsection*{L'ensemble des scindages, nouvel ensemble 3-3}

Les pages 7 et 8 reprennent sans la structure (A). Soit $S$ un ensemble 3-3 et
$\mathcal{G}$ son groupe d'automorphismes, extension
\[
1 \to \prod_{i \in T} \mathfrak{S}_{S_i} \to \mathcal{G} \to \mathfrak{S}_T \to 1 ,
\]
d'ordre $72$ : c'est le produit en couronne $\mathfrak{S}_3 \wr \mathbf{Z}/2\mathbf{Z}$. Soit
$S' \subset \mathcal{G}$ l'ensemble des éléments involutifs au-dessus de l'élément non
trivial de $\mathfrak{S}_T$, c'est-à-dire des scindages de l'extension. Pour $i \in T$
fixé, $S'$ est en bijection avec $\mathrm{Isom}(S_i, S_j)$, $j \neq i$, ou avec
l'ensemble des graphes $\Gamma \subset \prod_i S_i$ de ces bijections : il a six
éléments.

On a une surjection
\[
S' \longrightarrow T' := \bigwedge_{i \in T} \omega_{S_i},
\]
qui envoie une bijection sur l'identification des orientations qu'elle transporte,
et dont les fibres ont trois éléments. Ainsi $S'$ est à son tour un ensemble 3-3,
de classes indexées par $T'$. Un choix de (A) est un point de $T'$ ; au-dessus de
lui, la fibre est l'ensemble $I$ des trois antipodies, et l'autre fibre celui des
trois involutions renversantes, qu'on vient d'identifier à $\mathsf{T}$. Le groupe $G$
des pages 3 à 6 est le stabilisateur de ce point dans $\mathcal{G}$, et l'égalité
$G \simeq \mathfrak{S}_I \times \mathfrak{S}_{\mathsf{T}}$ dit que $G$ est le groupe des
automorphismes de l'ensemble 3-3 $S'$ qui préservent chacune de ses deux classes.

Enfin, un hexagone de triangles donnés équivaut à un appariement de ses sommets
opposés, donc à un point de $S'$ : la donnée d'une structure hexagonale sur $S$
équivaut à celle de la partition et d'un élément $s \in S'$\footnote{La page
poursuit « — celle d'un hexagone et d'un sommet dessus, à un » et s'arrête au
premier tiers du feuillet ; le sens de la fin de phrase n'est pas établi.}.

\pagerange{9}{14}
\section*{Plans sur $\mathbf{F}_3$ et carrés combinatoires (pages 9 à 14)}

\pagerange{9}{10}
\subsection*{Deux paires de droites, deux paires de quadrillages}

Soit $V$ un plan vectoriel sur $\mathbf{F}_3$. La droite projective $P(V)$ a quatre
points : $V$ contient quatre droites. Donnons-nous une paire $T$ de droites
$(D_i)_{i \in T}$, et soit $T' = P(V) \smallsetminus T$ la paire complémentaire. On a
\[
(1)\quad V = \bigoplus_{i \in T} D_i, \qquad (2)\quad V = \bigoplus_{i \in T'} D_i .
\]
Chaque $D_i^{*}$ a deux éléments, et l'on a des bijections canoniques
\[
(3) \qquad T' \simeq \bigwedge_{i \in T} D_i^{*}, \qquad T \simeq \bigwedge_{i \in T'} D_i^{*},
\]
la première induite par
\[
(4) \qquad (u_i)_{i \in T} \longmapsto \mathbf{F}_3 \cdot \Bigl(\sum_{i \in T} u_i\Bigr) :
\]
la somme d'un générateur de chaque droite de $T$ n'est sur aucune d'elles, donc
engendre une droite de $T'$ ; changer les deux signes ne change pas la droite, et
changer un seul signe fait passer à l'autre\footnote{Justification de l'édition ;
la page écrit (4) sans commentaire.}.

Soit $P$ un $V$-torseur, plan affine à neuf points. Pour $i \in P(V)$, notons
$\bar\imath$ l'autre élément de la paire, $T$ ou $T'$, qui contient $i$ : c'est
l'involution de $P(V)$ associée à la partition de type $(2,2)$ $\{T, T'\}$. Le
quotient
\[
(6) \qquad P_i = P/D_{\bar\imath} = P \wedge_V (V/D_{\bar\imath}), \qquad V/D_{\bar\imath} \simeq D_i ,
\]
est un $D_i$-torseur à trois éléments, et (1), (2) donnent deux décompositions
\[
(5) \qquad P \simeq \prod_{i \in T} P_i, \qquad P \simeq \prod_{i \in T'} P_i .
\]
Ainsi les neuf points de $P$ portent deux paires de partitions de type $(3,3,3)$
— les classes de parallélisme des quatre directions —, indexées l'une par $T$,
l'autre par $T'$\footnote{La page écrit « de type $(3,3)$ » pour la seconde paire.}.
Les deux partitions d'une même paire sont « orthogonales » : l'application
canonique de $P$ vers le produit des deux quotients est bijective, ce qui revient
à dire que chaque classe de l'une rencontre chaque classe de l'autre en exactement
un point — et il suffit pour cela que l'intersection soit toujours non vide, ou
toujours d'au plus un point.

\pagerange{11}{12}
\subsection*{L'équivalence entre les deux côtés}

Les décompositions (5) donnent des équivalences de groupoïdes de torseurs
\[
(7) \qquad \prod_{i \in T} \mathbf{Tors}(D_i) \ \simeq\ \mathbf{Tors}(V) \ \simeq\
\prod_{i \in T'} \mathbf{Tors}(D_i) .
\]
Pour une droite $D$ sur $\mathbf{F}_3$, $\mathbf{Tors}(D)$ s'identifie canoniquement au
groupoïde des ensembles $E$ à trois éléments munis d'une bijection
$\omega_E \simeq D^{*}$ : les deux générateurs de $D$ agissent sur un torseur par
les deux permutations circulaires, inverses l'une de l'autre. Par
$(P_i)_{i \in T} \mapsto \coprod_{i \in T} P_i$, le groupoïde
$\prod_{i \in T} \mathbf{Tors}(D_i)$ devient celui des ensembles 3-3 de classes
indexées par $T$ munis de bijections $\omega_{P_i} \simeq D_i^{*}$. Les deux
équivalences (7) fournissent donc une équivalence
\[
(8) \qquad
\left\{\begin{array}{l}
\text{ensembles 3-3 } (P_i)_{i \in T}, \\
\text{avec } D_i^{*} \simeq \omega_{P_i} \ (i \in T)
\end{array}\right\}
\ \xrightarrow{\ \simeq\ }\
\left\{\begin{array}{l}
\text{ensembles 3-3 } (P_j)_{j \in T'}, \\
\text{avec } D_j^{*} \simeq \omega_{P_j} \ (j \in T')
\end{array}\right\}
\]
Si l'on échange les rôles de $T$ et $T'$, le diagramme (7) est le même à symétrie
près, et l'équivalence obtenue est quasi-inverse de (8).

Sur les ensembles sous-jacents, (8) envoie $S = P_1 \sqcup P_2$ sur l'ensemble des
droites affines de $P = P_1 \times P_2$ dans les deux directions de $T'$. C'est,
comme on l'a dit dans « Le fil du dossier », l'ensemble $S'$ des scindages des
pages 7 et 8, et le point de $T'$ que choisit (A) à la page 3 est l'une de ces
deux directions\footnote{Ce rapprochement est de l'édition. L'indexation de (8),
où la classe des droites de direction $j$ est $P/D_j = P_{\bar\jmath}$, diffère de
celle de $S' \to T'$, qui envoie une droite sur sa direction, par l'involution
$j \mapsto \bar\jmath$.}.

\pagerange{12}{14}
\subsection*{Carrés combinatoires et carré dual}

Les deux groupoïdes de (8) ne dépendent que de la famille, indexée par
$P(V) = T \sqcup T'$, des ensembles à deux éléments $\omega_i = D_i^{*}$. Or une
famille $(\omega_i)_{i \in T}$ suffit à reconstituer, à isomorphisme près, les
droites $D_i = \omega_i \sqcup \{0\}$ — tout ensemble à deux éléments est de façon
unique un torseur sous $\mathbf{F}_3^{*} = \{\pm 1\}$ —, donc le plan
$V = \prod_{i \in T} D_i$ avec sa paire de droites, et la paire complémentaire.
D'autre part, une famille de deux ensembles à deux éléments est exactement la
donnée d'un \emph{carré combinatoire} $Q$ dont $T$ est l'ensemble des
diagonales : les sommets sont $\coprod_i \omega_i$, chaque sommet d'une diagonale
étant joint aux deux sommets de l'autre\footnote{« carré » est la lecture, dictée
par la suite, d'un mot souligné qu'on pourrait prendre pour « couvé ».}. D'où
une équivalence de groupoïdes
\[
(9) \qquad \{\text{carrés combinatoires } Q\} \ \simeq\
\{\text{plans sur } \mathbf{F}_3 \text{ munis d'une paire de droites}\},
\]
dans laquelle, en notant $d_i$ la diagonale $i$ de $Q$,
\[
(10) \qquad T \simeq \operatorname{diag} Q, \quad D_i^{*} \simeq d_i, \quad
S_Q = \coprod_{i \in T} D_i^{*}, \quad
A_Q = \prod_{i \in T} D_i^{*} \simeq V \smallsetminus \bigcup_{i \in T} D_i = \bigcup_{j \in T'} D_j^{*} .
\]
Les groupes d'automorphismes se correspondent : le groupe diédral d'ordre $8$ d'un
côté, les matrices monomiales de $GL_2(\mathbf{F}_3)$ de l'autre\footnote{Cette
vérification est de l'édition. La page écrit $V^{*} \smallsetminus \bigcup D_i^{*}$ pour
l'ensemble des vecteurs non nuls hors des deux droites.}.

La dernière égalité de (10) montre que le carré défini par la paire
complémentaire $(D_j)_{j \in T'}$ s'identifie au \emph{carré dual} $DQ$, dont les
sommets sont les arêtes de $Q$, deux arêtes étant adjacentes dans $DQ$ quand elles
ont un sommet commun : les arêtes $(u_1, u_2)$ et $(u_1, -u_2)$ ont pour sommes
deux vecteurs de directions différentes, et les arêtes $(u_1, u_2)$ et
$(-u_1, -u_2)$, sans sommet commun, deux vecteurs opposés d'une même droite. Si
$V(Q)$ désigne le plan associé à $Q$ par (9), on a donc un isomorphisme
\[
(11) \qquad \alpha_Q : V(DQ) \xrightarrow{\ \sim\ } V(Q), \qquad e = (u_1, u_2) \longmapsto u_1 + u_2 ,
\]
et, symétriquement, en identifiant $DDQ$ à $Q$ — un sommet de $Q$ est l'arête de
$DQ$ formée des deux arêtes de $Q$ qui s'y rencontrent —,
\[
(12) \qquad \alpha_{DQ} : V(Q) = V(DDQ) \xrightarrow{\ \sim\ } V(DQ) .
\]
La page numérote cette dernière formule (11), comme la précédente\footnote{La
suite, (13), demande (12).}.

\textbf{La relation (13).}
\[
\alpha_Q\, \alpha_{DQ} = -\mathrm{id}_{V(Q)} = V(a_Q),
\qquad \alpha_{DQ}\, \alpha_Q = -\mathrm{id}_{V(DQ)} = V(a_{DQ}),
\]
où $a_Q$ est l'antipodie de $Q$, l'automorphisme qui échange les deux sommets de
chaque diagonale. La page l'énonce ; en voici la raison\footnote{Vérification de
l'édition.}. Un sommet $u_1 \in d_1$ de $Q$ est commun aux arêtes $e = (u_1, u_2)$
et $e' = (u_1, -u_2)$, qui sont sur les deux diagonales de $DQ$ ; donc
$\alpha_{DQ}(u_1) = e + e'$, et
\[
\alpha_Q(e + e') = (u_1 + u_2) + (u_1 - u_2) = 2u_1 = -u_1 ,
\]
puisque $2 = -1$ dans $\mathbf{F}_3$. De même pour les sommets de $d_2$. La relation
est ainsi propre à la caractéristique $3$ : passer du carré à son dual et revenir
ne donne pas l'identité, mais l'antipodie.

\pagerange{15}{18}
\section*{Reconstituer un module libre de rang $2$ à partir de ses droites (pages 15 à 18)}

La page 13 reconstituait un plan sur $\mathbf{F}_3$ à partir d'une paire de droites
données. Les pages 15 à 18 reprennent la question sans choisir de paire, et sur
un anneau quelconque\footnote{La page commence par « un Module loc. libre » sur
« un topos loc. annelé », biffés au profit d'un module libre de rang $2$ sur un
anneau.}.

Soit $\Lambda$ un anneau commutatif et $V$ un $\Lambda$-module libre de rang $2$.
On note
\begin{itemize}
  \item $V^{*}$ l'ensemble des vecteurs \emph{unimodulaires} de $V$, ceux qui
    font partie d'une base, et $\Lambda^{*}$ le groupe des unités ;
  \item $P(V)^{!} = V^{*}/\Lambda^{*}$, l'ensemble des droites de $V$ — facteurs
    directs de rang $1$ — qui sont libres, partie de la droite projective $P(V)$ ;
  \item pour $x \in P(V)^{!}$, $D_x$ la droite correspondante et $T_x = D_x^{*}$ le
    $\Lambda^{*}$-torseur de ses générateurs ;
  \item $\omega_V = \det(V)^{*}$, le $\Lambda^{*}$-torseur des générateurs de
    $\det V = \Lambda^2 V$.
\end{itemize}
Deux droites $x, y$ sont \emph{disjointes} si $V = D_x \oplus D_y$ : sur un corps,
cela veut dire $x \neq y$ ; sur un anneau local, que leurs réductions modulo
l'idéal maximal sont distinctes. C'est le cas si et seulement si, pour
$u \in T_x$ et $v \in T_y$, $u \wedge v$ engendre $\det V$ ; l'application
$(u, v) \mapsto u \wedge v$ passe alors au produit contracté et donne un isomorphisme
de $\Lambda^{*}$-torseurs
\[
(5) \qquad T_x \wedge_{\Lambda^{*}} T_y \xrightarrow{\ \sim\ } \omega_V ,
\]
avec $u \wedge v = -\,v \wedge u$, où $-u = (-1)\cdot u$.

\textbf{La relation fondamentale.} Soient $x, y$ disjointes, $z$ quelconque, et
$u \in T_x$, $v \in T_y$, $w \in T_z$. Alors
\[
u \wedge v = v \wedge w = w \wedge u \iff u + v + w = 0 ,
\]
et ces conditions entraînent que $z$ est disjointe de $x$ et de $y$. En effet,
$(u, v)$ est une base ; écrivons $w = au + bv$. Alors $v \wedge w = -a\,(u \wedge v)$
et $w \wedge u = -b\,(u \wedge v)$, et les égalités demandées équivalent à
$a = b = -1$\footnote{Démonstration de l'édition ; la page dit « on vérifie de
suite ». L'énoncé vaut sur tout anneau.}. Sur $\mathbf{F}_3$, c'est la
bijection (4) de la page 9 : pour $x, y$ données, la droite $z$ de $w = -(u + v)$
ne dépend que de la classe de $(u, v)$ dans $T_x \wedge T_y$.

\textbf{La reconstruction.} Grothendieck affirme que, du moins si $\Lambda$ est un
corps, $V$ se reconstitue à partir des données suivantes :
\begin{itemize}
  \item[a)] l'ensemble $P = P(V)^{!}$ ;
  \item[b)] la famille des $\Lambda^{*}$-torseurs $(T_x)_{x \in P}$, i.e. le torseur
    relatif $\coprod_x T_x = V^{*}$ sur $P$ — ce qui revient, avec a), à se donner
    l'ensemble $V^{*}$ muni de l'action libre de $\Lambda^{*}$ ;
  \item[c)] le $\Lambda^{*}$-torseur $\omega = \omega_V$ ;
  \item[d)] la relation de disjonction sur $P$, ou sur $V^{*}$, où elle est stable
    par $\Lambda^{*}$ en chaque variable — triviale si $\Lambda$ est un corps ;
  \item[e)] pour $x, y$ disjointes, les isomorphismes (5).
\end{itemize}
Sur un corps, on pose $V = V^{*} \sqcup \{0\}$ et $D_x = T_x \sqcup \{0\}$, qui est une
droite sur $\Lambda$ ; on définit $u + v$ quand l'un est nul, quand $u$ et $v$ sont
sur une même droite (par l'addition de $D_x$), et, quand ils sont disjoints, par
$u + v = -w$, où $w$ est l'unique élément de $V^{*}$ disjoint de $u$ et $v$ tel que
$u \wedge v = v \wedge w = w \wedge u$. Pour que cette construction réussisse, il
faut des propriétés des données a) à e), dont la page n'écrit que la première :
\begin{itemize}
  \item[a)] pour $x, y, z$ deux à deux disjointes, il existe $u \in T_x$, $v \in
    T_y$, $w \in T_z$ avec $u \wedge v = v \wedge w = w \wedge u$, et l'on peut même
    imposer l'un des trois.
\end{itemize}
Cette propriété est vraie, sur tout anneau : si $w$ engendre $D_z$, sa
décomposition $-w = u + v$ dans $D_x \oplus D_y$ a des composantes qui engendrent
$D_x$ et $D_y$, puisque $z$ est disjointe de chacune\footnote{Justification de
l'édition. Les autres propriétés, qui doivent assurer que l'addition ainsi
définie est associative et que $V$ est un espace vectoriel, ne sont pas écrites
: la page s'arrête au tiers inférieur du feuillet.}.

\pagerange{19}{20}
\section*{Polygones combinatoires (pages 19 et 20)}

Ces deux pages, d'une main plus rapide, posent les définitions dont les
hexagones de la page 2 et les carrés de la page 13 sont des cas. Le foncteur
qui associe à un polygone topologique sa structure combinatoire est
essentiellement surjectif : tout polygone combinatoire se réalise
topologiquement\footnote{La page écrit « Pol top $\to$ Pol comb est ess.
surjectif » et, en interligne, « Foncteur réalisation topologique », lecture
incertaine ; la suite de la ligne est illisible.}.

\textbf{Trois points de vue.} Un polygone combinatoire $\Pi$ peut se définir :
\begin{enumerate}
  \item comme un ensemble ordonné $S \sqcup A$, où $S$ (les sommets) est l'ensemble
    des éléments minimaux et $A$ (les arêtes) celui des éléments maximaux, chaque
    sommet étant strictement majoré par exactement deux arêtes et chaque arête
    strictement minorée par exactement deux sommets ; la dualité échange $S$ et
    $A$ en renversant l'ordre ;
  \item par l'ensemble de base $S$ et une partie $A \subset \mathfrak{P}_2(S)$ de
    parties à deux éléments, chaque sommet appartenant à exactement deux arêtes —
    point de vue valable si les composantes connexes ont au moins trois sommets,
    le digone ayant deux arêtes de mêmes extrémités ;
  \item pour un polygone connexe d'au moins trois sommets, par une paire d'ordres
    circulaires sur $S$ (ou sur $A$) inverses l'un de l'autre, c'est-à-dire de
    permutations circulaires $f, f^{-1}$.
\end{enumerate}
Le premier point de vue décrit les graphes bipartis $2$-réguliers, c'est-à-dire
les réunions disjointes de cycles, digones et cycles infinis compris\footnote{Cette identification
est de l'édition.}.

\textbf{Orientations.} Orienter une arête, c'est choisir son origine ; son bord est
alors « extrémité moins origine ». Une orientation de $\Pi$ est un choix
d'orientation de chaque arête tel qu'en chaque sommet $s$ les deux orientations
induites soient opposées : une arête arrive en $s$, l'autre en part. Dans le
troisième point de vue, c'est le choix de l'un des deux ordres circulaires.

\textbf{Théorème.} Soit $\Pi$ un polygone combinatoire connexe, d'au moins trois
sommets\footnote{La page ne suppose rien pour a) et b) et ajoute seulement, en
NB, l'hypothèse « comp. d'ordre $\geqslant 3$ ». Sans connexité, $\Pi$ a
$2^{c}$ orientations, $c$ étant le nombre de composantes. Pour le digone, a) et
b) valent encore, mais c) est faux : les deux orientations donnent le même
automorphisme.}.
\begin{enumerate}
  \item[a)] $\Pi$ admet exactement deux orientations.
  \item[b)] À toute orientation $\omega$ correspond un unique automorphisme $f_\omega$
    de $\Pi$ qui respecte $\omega$ et envoie chaque sommet $s$ sur l'extrémité de
    l'arête d'origine $s$ ; il envoie chaque arête sur l'arête suivante.
  \item[c)] $\omega \mapsto f_\omega$ est une bijection de l'ensemble des orientations
    sur l'ensemble des automorphismes $f$ de $\Pi$ tels que, pour tout sommet $s$,
    $f(s)$ soit adjacent à $s$, et, pour toute arête $a$, $f(a) \neq a$.
\end{enumerate}
Pour a), l'orientation d'une arête se propage de proche en proche le long du
cycle. Pour c), la condition sur les arêtes n'est nécessaire que pour le carré,
dont les deux réflexions par rapport aux médiatrices de côtés envoient elles
aussi chaque sommet sur un voisin, mais fixent deux arêtes\footnote{La page
formule c) par référence à la condition 2) de b), qui fait intervenir l'origine
et donc l'orientation elle-même. La formulation retenue est celle de la marge
gauche, « $f(s)$ soit adjacent à $s$ ; $f(a) \neq a$ », dont plusieurs mots sont
illisibles. Les automorphismes décrits sont les deux rotations d'un cran.}.
Autrement dit : une orientation d'un polygone, c'est la même chose qu'une rotation
élémentaire. La page se termine sur la remarque qu'orienter $\Pi$ en un sommet $s$
revient à choisir, parmi les deux arêtes en $s$, l'arête « origine » et l'arête
« extrémité », et annonce la comparaison avec l'orientation d'un segment et d'une
variété topologique de dimension $1$, sans la faire.

\end{document}
